% SPDX-License-Identifier: CC-BY-SA-4.0
% Author: Matthieu Perrin

\begingroup

\begin{exercice}[Équivalence d'expressions rationnelles]\label{exo:lexing/arden/simplify_regexp}

  \begin{question}
  \item Montrer que $a(ba)^\star  b$ est solution de $X = ab X \mid ab$. 
  \item En déduire $a(ba)^\star  b \equiv (ab)^+$.
  \item Montrer que $a(a\mid ba)^\star$ est solution de $X = (a \mid ab) X \mid a$.
  \item En déduire $a(a\mid ba)^\star \equiv (a\mid ab)^\star  a$
  \end{question}

\end{exercice}

\begin{correction}

  \begin{question}
  \item $\begin{array}[t]{rcl}
    ab\, \underline{a(ba)^\star  b} \mid ab
    &=& a  (ba(ba)^\star \mid \varepsilon ) b \\
    &=& \underline{a  (ba)^\star b} \\
  \end{array}$
  \item D'après le lemme d'Arden, l'unique solution de $X = ab X \mid ab$ est $(ab)^+$, donc $a(ba)^\star  b \equiv (ab)^+$. 
  \item $\begin{array}[t]{rcl}
    (a \mid ab) \underline{a(a\mid ba)^\star} \mid a
    &=& a \,\underline{a(a\mid ba)^\star} \mid ab \,\underline{a(a\mid ba)^\star} \mid a\\
    &=&  a (a(a\mid ba)^\star \mid ba(a\mid ba)^\star \mid \varepsilon)\\
    &=& \underline{a(a\mid ba)^\star} \\
  \end{array}$
  \item D'après le lemme d'Arden, l'unique solution de $X = (a\mid ab) X \mid a$ est $(a\mid ab)^\star a$, \\donc $a(a\mid ba)^\star \equiv (a\mid ab)^\star  a$. 
  \end{question}

\end{correction}

\endgroup
\endinput
